\chapter{G\'en\'erateur de Dataset Synth\'etique} \label{ch_dataset}

Cet outil g\'en\`ere un dataset synth\'etique d'images en noir et blanc contenant des formes g\'eom\'etriques (cercles et triangles \'equilat\'eraux) pour des t\^aches de vision par ordinateur. L'objectif est de d\'etecter si un triangle est \`a l'int\'erieur ou \`a l'ext\'erieur d'un cercle.

\section{Fonctionnalit\'es}

\begin{itemize}
    \item \textbf{G\'en\'eration Synth\'etique} : Cr\'ee des images de 128x128 pixels avec un cercle et un triangle \'equilat\'eral.
    \item \textbf{Classes} :
    \begin{itemize}
        \item \texttt{inside} : Le triangle est strictement contenu \`a l'int\'erieur du cercle.
        \item \texttt{outside} : Le triangle est strictement disjoint du cercle.
    \end{itemize}
    \item \textbf{Style Visuel} :
    \begin{itemize}
        \item Fond : Blanc
        \item Cercle : Gris rempli
        \item Triangle : Noir rempli
    \end{itemize}
    \item \textbf{Visualisation} : Inclut un outil pour inspecter le dataset g\'en\'er\'e dans un navigateur web.
    \item \textbf{Param\'etrable} : Hautement configurable via \texttt{config.py}.
\end{itemize}

\section{Installation}

\begin{enumerate}
    \item Assurez-vous d'avoir Python 3 install\'e.
    \item Installez les d\'ependances requises :
\end{enumerate}

\begin{verbatim}
pip install Pillow
\end{verbatim}

\section{Utilisation}

\subsection{1. Configurer le Dataset}
\'Editez \texttt{config.py} pour ajuster les param\`etres :
\begin{itemize}
    \item \texttt{NUM\_IMAGES} : Nombre total d'images \`a g\'en\'erer.
    \item \texttt{IMAGE\_SIZE} : R\'esolution (par d\'efaut 128x128).
    \item \texttt{MIN\_CIRCLE\_RADIUS} / \texttt{MAX\_CIRCLE\_RADIUS} : Contraintes de taille pour les cercles.
    \item \texttt{MIN\_TRIANGLE\_SIDE} / \texttt{MAX\_TRIANGLE\_SIDE} : Contraintes de taille pour les triangles.
    \item \texttt{RATIO\_INSIDE} : Proportion d'\'echantillons "inside" vs "outside".
\end{itemize}

\subsection{2. G\'en\'erer le Dataset}
Ex\'ecutez le script de g\'en\'eration :

\begin{verbatim}
python3 dataset_generator.py
\end{verbatim}

Cela cr\'eera un r\'epertoire \texttt{generated\_dataset/} contenant :
\begin{itemize}
    \item \texttt{images/} : Les images PNG g\'en\'er\'ees.
    \item \texttt{annotations.csv} : \'Etiquettes pour chaque image.
    \item \texttt{annotations.json} : \'Etiquettes au format JSON.
    \item \texttt{index.html} : Une galerie pour visualiser les images.
\end{itemize}

\subsection{3. Visualiser les R\'esultats}
Pour voir le dataset, ex\'ecutez :

\begin{verbatim}
python3 visualizer.py
\end{verbatim}

Cela v\'erifiera la sortie et ouvrira automatiquement \texttt{generated\_dataset/index.html} dans votre navigateur par d\'efaut, vous permettant de filtrer et parcourir les images.
